\documentclass[a4paper, serif, 10pt]{moderncv}

%% moderncv
% style and color
\moderncvstyle{banking}
\moderncvcolor{black}

% renew moderncv command to adapt for CJK colon
\renewcommand*{\cvitem}[3][.25em]{%
  \ifstrempty{#2}{}{\hintstyle{#2}：}{#3}%
  \par\addvspace{#1}}

\renewcommand*{\cvitemwithcomment}[4][.25em]{%
  \savebox{\cvitemwithcommentmainbox}{\ifstrempty{#2}{}{\hintstyle{#2}：}#3}%
  \setlength{\cvitemwithcommentmainlength}{\widthof{\usebox{\cvitemwithcommentmainbox}}}%
  \setlength{\cvitemwithcommentcommentlength}{\maincolumnwidth-\separatorcolumnwidth-\cvitemwithcommentmainlength}%
  \begin{minipage}[t]{\cvitemwithcommentmainlength}\usebox{\cvitemwithcommentmainbox}\end{minipage}%
  \hfill% fill of \separatorcolumnwidth
  \begin{minipage}[t]{\cvitemwithcommentcommentlength}\raggedleft\small\itshape#4\end{minipage}%
  \par\addvspace{#1}}

%% page layout/margins
\usepackage[top=2.5cm, bottom=2.5cm, left=1.5cm, right=1.5cm]{geometry}

\nopagenumbers{}

%% Babel config for English language
\usepackage[english]{babel}

%% fontspec
\usepackage{fontspec}

\IfFontExistsTF{Linux Libertine}{
  \setmainfont[Ligatures={TeX, Common}, Numbers=Lining, ItalicFont=Linux Libertine]{Linux Libertine}
}{}
\IfFontExistsTF{Linux Libertine O}{
  \setmainfont[Ligatures={TeX, Common}, Numbers=Lining, ItalicFont=Linux Libertine O]{Linux Libertine O}
}{}

%% CTeX
% CJK support, used to show CJK characters in the resume
%
% - fontset=none: disable builtin fontset but instead set the CJK font manually
% - heading=false: disable ctex heading
% - punct=kaiming: use kaiming punctuations styles for CJK
% - scheme=plain: use plain scheme, do not override \normalsize font size
% - space=auto: space settings for CJK characters
%
% ref:
% - http://ctan.mirrorcatalogs.com/language/chinese/ctex/ctex.pdf
\usepackage[UTF8, heading=false, punct=kaiming, scheme=plain, space=auto]{ctex}

\IfFontExistsTF{Noto Serif CJK SC}{
  \setCJKmainfont{Noto Serif CJK SC}
}{}
\IfFontExistsTF{Noto Sans CJK SC}{
  \setCJKsansfont{Noto Sans CJK SC}
}{}

%% line spacing
\usepackage{setspace}
\setstretch{1.125}

%% URL styling - use normal text instead of monospace
\urlstyle{same}

% Auto-underline all links
% Defer until \begin{document} so hyperref is already loaded
\AtBeginDocument{%
  \let\oldhref\href
  \renewcommand{\href}[2]{\oldhref{#1}{\underline{#2}}}%
}

%% Patch \httplink and \httpslink to handle full URLs
% The original moderncv macros always prepend http:// or https:// to URLs.
% This causes issues when the URL already has a protocol (e.g., https://example.com)
% resulting in malformed URLs like https://https://example.com.
% This patch checks if the URL already contains "://" and uses it directly if so.
% Note: We use \str_set:Nx (x-expansion) to expand macros like \@homepage first.
\ExplSyntaxOn
\str_new:N \g_yamlresume_url_str

\RenewDocumentCommand{\httpslink}{O{}m}{
  \str_set:Nx \g_yamlresume_url_str {#2}
  \str_if_in:NnTF \g_yamlresume_url_str {://}
    { \url{#2} }
    { \url{https://#2} }
}

\RenewDocumentCommand{\httplink}{O{}m}{
  \str_set:Nx \g_yamlresume_url_str {#2}
  \str_if_in:NnTF \g_yamlresume_url_str {://}
    { \url{#2} }
    { \url{http://#2} }
}
\ExplSyntaxOff

\name{陳慧敏}{}
\title{數據分析師 / 商業智能分析}
\phone[mobile]{+852 9123 4567}
\email{waimanch.chan@outlook.com}
\homepage{https://wmchan-data.github.io}

\address{中國香港，香港島，香港，香港中環金融街8號}{}{}

\extrainfo{{\small \faLinkedin}\ \href{https://www.linkedin.com/in/wmchan-data}{@wmchan-data} {} {} {} • {} {} {}
{\small \faGithub}\ \href{https://github.com/wmchan-data}{@wmchan-data}}

\begin{document}

\maketitle

\section{簡介}

\cvline{}{\begin{itemize}
\item 具 5 年商業數據分析經驗，擅長從龐大數據中發掘業務洞察，支援管理層決策
\item 精通 SQL、Python 及 Power BI，熟悉 ETL 流程、數據倉庫及 A/B 測試設計
\item 曾主導多個跨部門數據項目，包括客戶流失預測、銷售預測及風險監察儀表板
\item 善於以視覺化方式向非技術持分者清晰傳遞分析結果，推動數據驅動文化
\end{itemize}}
\section{教育背景}

\cventry{2017年9月–2021年7月}
        {學士，數據科學及商業智能，成績：3.6 / 4.0}
        {香港科技大學}
        {\url{https://www.hkust.edu.hk}}
        {}
        {\begin{itemize}
\item 主修課程涵蓋統計模型、機器學習、數據庫管理及商業智能
\item 畢業專題為「香港中小企客戶流失預測模型」，獲學系最佳專題獎
\item 副修計算機科學，具備扎實的程式設計及數據處理基礎
\end{itemize}
\textbf{課程}：數據庫管理系統、統計推斷、機器學習、商業分析、資料視覺化、時間序列分析、風險管理與金融科技、大數據計算}
\section{工作經歷}

\cventry{2023年7月至今}
        {數據分析師}
        {香港金融管理局}
        {\url{https://www.hkma.gov.hk}}
        {}
        {\begin{itemize}
\item 為貨幣穩定及金融體系監察團隊建立自動化數據儀表板，每日處理超過 500 萬筆交易數據
\item 利用 Python 及 SQL 整合內部數據倉庫，縮短報表編製時間約 40\%
\item 與經濟研究部門合作，以時間序列模型分析銀行體系流動性風險
\item 定期向高級管理層提交數據分析報告，協助制訂監管政策
\end{itemize}
\textbf{關鍵字}：Python、SQL、Tableau、時間序列分析、風險監察}

\cventry{2021年8月–2023年6月}
        {商業數據分析助理}
        {恒生銀行}
        {\url{https://www.hangseng.com}}
        {}
        {\begin{itemize}
\item 協助零售銀行部門設計客戶消費行為分析框架，支援信用卡及貸款產品策略
\item 建立 Power BI 儀表板追蹤客戶忠誠度及產品滲透率，每月節省約 20 工時
\item 進行 A/B 測試評估數碼推廣活動成效，並提出具體優化建議
\item 整理及清洗來自多個來源的客戶數據，確保數據品質及合規
\end{itemize}
\textbf{關鍵字}：Power BI、A/B 測試、數據清洗、客戶分析}
\section{語言}

\cvline{中文}{母語或雙語水平 \hfill \textbf{關鍵字}：粵語、普通話、繁體中文}
\cvline{英語}{完全專業水平 \hfill \textbf{關鍵字}：IELTS 7.5、職場英語、技術簡報}
\section{專業技能}

\cvline{數據分析}{專家 \hfill \textbf{關鍵字}：Python、Pandas、NumPy、Jupyter Notebook、統計分析}
\cvline{資料庫及 SQL}{專家 \hfill \textbf{關鍵字}：PostgreSQL、MySQL、BigQuery、SQL Server、資料倉庫}
\cvline{商業智能及視覺化}{高級 \hfill \textbf{關鍵字}：Power BI、Tableau、Looker Studio、Excel、儀表板設計}
\cvline{機器學習基礎}{中級 \hfill \textbf{關鍵字}：Scikit-learn、迴歸分析、分類模型、集群分析}
\cvline{數據工程}{中級 \hfill \textbf{關鍵字}：ETL、Apache Airflow、dbt、資料建模}
\section{榮譽}

\cventry{2021年6月}
        {香港科技大學}
        {香港科技大學院長嘉許名單}
        {}
        {}
        {連續三個學期入選院長嘉許名單，表揚學業成績優異及對大學社區的貢獻。}
\section{證書}

\cventry{2024年1月}
        {Microsoft}
        {Microsoft 認證：Power BI 數據分析師助理（PL-300）}
        {\url{https://learn.microsoft.com/en-us/certifications/power-bi-data-analyst-associate/}}
        {}
        {}

\cventry{2023年3月}
        {Tableau}
        {Tableau Desktop 專家認證}
        {\url{https://www.tableau.com/learn/certification/desktop-specialist}}
        {}
        {}

\cventry{2022年8月}
        {Google}
        {Google 數據分析專業證書}
        {\url{https://www.coursera.org/professional-certificates/google-data-analytics}}
        {}
        {}
\section{出版刊物}

\cventry{2023年10月}
        {香港虛擬銀行用戶行為研究報告}
        {香港金融科技協會}
        {\url{https://www.hkfta.org/research/virtual-bank-users-2023}}
        {}
        {\begin{itemize}
\item 以 1,200 名受訪者數據分析虛擬銀行用戶的採用障礙及使用習慣
\item 提出改善用戶體驗及數碼營銷策略的具體建議
\end{itemize}}
\section{項目}

\cventry{2023年3月–2023年6月}
        {一個以 Python 及 Power BI 開發的互動式零售銷售預測平台，整合公開數據與 API，為管理層提供決策支援。}
        {香港零售銷售預測儀表板}
        {\url{https://github.com/wmchan-data/retail-forecast-dashboard}}
        {}
        {\begin{itemize}
\item 利用 Python 及 Prophet 建立月度銷售預測模型，平均絕對百分比誤差低於 12\%
\item 以 Power BI 建立互動儀表板，讓管理層掌握不同店舖及產品類別的銷售趨勢
\item 設計自動化數據更新流程，每日從 API 抽取數據並重新訓練模型
\end{itemize}
\textbf{關鍵字}：Python、Prophet、Power BI、時間序列、API}
\section{興趣愛好}

\cvline{開放資料}{香港政府資料一線通、資料新聞學、開放街圖}
\cvline{運動}{跑步、行山、瑜伽}
\section{志願服務}

\cventry{2023年1月至今}
        {數據分析義工}
        {香港紅十字會}
        {\url{https://www.redcross.org.hk}}
        {}
        {\begin{itemize}
\item 為籌款活動整理及分析歷年捐贈者數據，識別高潛在捐款群組
\item 建立活動成效儀表板，協助團隊更有效地分配宣傳資源
\end{itemize}}

\end{document}